\scanpage{100}
% archive sheet number: 99
\uline{2. Formalism of Algebraic Cohomology Classes}.

\uline{Cohomology Class Associated to a Cycle}. Let $C^*(X)$ be the intersection\inserted{\handwritten{ring}} of cycles on $X$,\inserted{\handwritten{modulo rational equivalence \unsure{(the Chow ring of $(X)$|(the Chow ring of $(K)$}}} tensored by $Q$. Then, for each $i\handwritten{\leq}n$, we define a homomorphism
\[
\handwritten{\rho_X:C^i(X)\longrightarrow H_\ell^{2i}(X)}
\]
as follows. Let $Z$ be an irreducible subscheme of codimension $i$. The l-adic Poincare duality theorem asserts that the cup-product pairing
\[
\handwritten{H_\ell^j(X)\times H_\ell^{2n-j}(X)\longrightarrow H_\ell^{2n}(X)\overset{\unsure{\kappa_X|K_X}}{\overset{\sim}{\longrightarrow}}Q_\ell\quad(j=0,\ldots,2n)}
\]
is a perfect duality of finite dimensional vector spaces, and so the element $\handwritten{\rho_X(Z)\in H_\ell^{2i}(X)}$ is uniquely determined by the condition
\[
\handwritten{\kappa_X\bigl(\alpha\cdot\rho_X(Z)\bigr)=\kappa_Z\bigl(\operatorname{res}_Z(\alpha)\bigr)\quad\text{for all }\alpha\in H_\ell^{2n-2i}(X).}
\]

We then extend $\handwritten{\rho_X}$ to all of $C^{\handwritten{i}}(X)$ by linearity. Further, this homomorphism takes intersection product into cup-product, and thus it extends to a homomorphism of \handwritten{``}algebras\handwritten{''} \deleted{on}\inserted{\handwritten{over}} $Q$\unsure{.}
\[
\handwritten{\rho_X:C^*(X)\longrightarrow H_\ell^*(X)..}
\]

The image of a cycle under this homomorphism .is called the \uline{cohomology class} associated with \unsure{thet|that} cycle, and we denote $\handwritten{\rho_X}(C^*(X))$ by $C_{\handwritten{\ell}}^*(X)$. A cycle is defined to be $Q_{\handwritten{\ell}}$-cohomologically equivalent to $0$ if it is in the kernel of the homomorphism $\handwritten{\rho_X.}$
